%% Lebenslauf-Grundgeruest - [YOUR_NAME]
%%
%% Das ist eine kompakte Ausgangsvorlage - KEINE vollstaendige Sammlung aller
%% Kompetenzen und Stationen. Eine Seite ist das bevorzugte Ziel; eine zweite
%% Seite ist erlaubt, wenn sie relevante, nicht redundante Belege bewahrt.
%%
%% Datenquelle fuer jede inhaltliche Aussage ist
%% .claude/skills/job-application-assistant/01-candidate-profile.md.
%% Beim Zuschneiden auf eine Stelle wird von dort ausgewaehlt und hier
%% eingesetzt - nicht umgekehrt. Diese Datei traegt die Struktur und die
%% ATS-Korrekturen im Praeambel-Block, nicht den Datenbestand.
%%
%% SETUP: /setup ausfuehren, um echte Daten einzutragen.
%% Kompilieren mit: lualatex main_example.tex
%% (pdflatex scheitert auf modernem MiKTeX oft an fontawesome5.)
%%
%% SEITENBUDGET: bevorzugt 1 Seite, harte Obergrenze 2 Seiten. Vor einer
%% zweiten Seite Redundanz entfernen; relevante Belege nicht nur fuer eine
%% kuenstliche Ein-Seiten-Regel streichen. Mehr als 2 Seiten ist ein Fehler.
%%
%% PLATZHALTER-SYNTAX: Platzhalter am Anfang eines \item immer klammern -
%% \item {[Text]}, niemals \item [Text]. Ohne Klammern liest LaTeX die
%% eckigen Klammern als optionales Label-Argument von \item; der Text
%% landet im linken Rand und wird an der Seitenkante abgeschnitten.

\documentclass[11pt,a4paper,sans]{moderncv}
\moderncvstyle{banking}
\moderncvcolor{blue}

% Vor- und Nachname sowie Abschnittsueberschriften in moderncv-Blau (color1).
% Banking laesst diese unter lualatex+MiKTeX sonst schwarz.
\renewcommand*{\firstnamestyle}[1]{{\fontsize{34}{36}\bfseries\upshape\color{color1}#1}}
\renewcommand*{\lastnamestyle}[1]{{\fontsize{34}{36}\bfseries\upshape\color{color1}#1}}
\renewcommand*{\sectionstyle}[1]{{\sectionfont\color{color1}#1}}

% ATS: Die FontAwesome-Kontaktsymbole haben kein Unicode-Mapping und werden
% von pdftotext als Glyphenname ("MOBILE-ANDROID-ALT") oder U+FFFD extrahiert.
% Leeren - Adresse, Nummer und E-Mail stehen ohnehin als Text daneben.
\renewcommand*{\mobilephonesymbol}{}
\renewcommand*{\phonesymbol}{}
\renewcommand*{\emailsymbol}{}
\renewcommand*{\addresssymbol}{}
\renewcommand*{\homepagesymbol}{}

% ATS: Die Standard-Aufzaehlungsmarke von moderncv ist ein kleiner Kreis ohne
% Unicode-Mapping und erscheint als Fremdzeichen in jeder Bullet-Zeile.
\renewcommand*{\labelitemi}{-}

% ATS: Das Trennzeichen der Kontaktzeile ist \rmfamily\textbullet und landet
% als blosses 0xB7-Byte im Textlayer - ungueltiges UTF-8, wird zu U+FFFD.
\renewcommand*{\makeheaddetailssymbol}{\hspace{0.8em}\textbar\hspace{0.8em}}

\usepackage[ngerman]{babel}
\usepackage{hyperref}
\hypersetup{
    colorlinks=true,
    linkcolor=blue,
    filecolor=magenta,
    urlcolor=blue,
    pdftitle={[YOUR_NAME] - Lebenslauf},
    pdfauthor={[YOUR_NAME]},
    pdfsubject={Lebenslauf},
}
\usepackage[scale=0.80]{geometry}
\usepackage{needspace}

% ATS: Blocksatz laesst LaTeX trennen und zerreisst deutsche Komposita ueber
% den Zeilenumbruch ("Investi-\ntionen"), sodass ein Screening-Parser das
% Schluesselwort nie findet. Flattersatz plus prohibitive Trennstrafe haelt
% jeden Begriff zusammen - auch echte Bindestriche wie Soll-Ist-Vergleich.
\raggedright
\hyphenpenalty=10000
\exhyphenpenalty=10000

% ------------------------------------------------------------
%     EINTRAGSFORMAT: EINSPALTIG (ATS)
% ------------------------------------------------------------
% Das zweispaltige \cventry von banking setzt Ort und Datum in eine
% rechtsbuendige Spalte. Im Rohtext-Modus von pdftotext erscheinen diese
% Felder als eigener Block hinter allen Eintraegen, sodass ein Parser
% Datumsangaben dem falschen oder keinem Arbeitgeber zuordnet.
% Einspaltig extrahiert jeder Eintrag sequenziell und vollstaendig.
%
% Argumente: {1 Datum}{2 Titel}{3 Organisation}{4 Ort}{5 Note}{6 Beschreibung}
\renewcommand*{\cventry}[6]{%
  \needspace{4\baselineskip}%
  \noindent{\bfseries #3}%
  \ifx&#2&\else, #2\fi\par
  \noindent{\itshape #1%
    \ifx&#4&\else\ \textbar\ #4\fi
    \ifx&#5&\else\ \textbar\ #5\fi}\par
  \vspace{1pt}#6\par\vspace{5pt}}

% ------------------------------------------------------------
%     ABSCHNITTSNAMEN
% ------------------------------------------------------------
% Lokalisierung ausschliesslich ueber diesen Block. Abschnittsueberschriften
% uebersetzen sich nicht von selbst - siehe 05-cv-templates.md.

% DEUTSCH (Standard):
\newcommand{\headingCompetencies}{Kernkompetenzen}
\newcommand{\headingExperience}{Berufserfahrung}
\newcommand{\headingProjects}{Projekte}
\newcommand{\headingEducation}{Ausbildung}
\newcommand{\headingLanguages}{Sprachen}
\newcommand{\headingAwards}{Auszeichnungen}

% ENGLISCH - bei englischsprachiger Ausschreibung diesen Block aktivieren und
% den deutschen auskommentieren. Zusaetzlich zu aendern:
%   - \usepackage[ngerman]{babel} -> \usepackage[english]{babel}
%   - \glqq ... \grqq{} im Kurzprofil -> "..." (\glqq stammt aus babel-ngerman
%     und ist unter [english] nicht definiert, der Lauf bricht sonst ab)
%   - Datumsformat und Notenangaben ("Note: 1,7" -> "Grade: 1.7 (German scale)")
%
% \newcommand{\headingCompetencies}{Core Competencies}
% \newcommand{\headingExperience}{Professional Experience}
% \newcommand{\headingProjects}{Projects}
% \newcommand{\headingEducation}{Education}
% \newcommand{\headingLanguages}{Languages}
% \newcommand{\headingAwards}{Honors and Awards}

% Persoenliche Daten
\name{[Vorname]}{[Nachname]}
\address{[Strasse Hausnummer, PLZ Ort]}{}{}
% Die Platzhalternummer hat absichtlich die Form einer echten Nummer: die
% Pruefung "phone in text" in tools/verify_pdf.py sucht eine Ziffernfolge,
% und ein Platzhalter aus X-en wuerde die eigene Beispieldatei durchfallen
% lassen. /setup ersetzt den Wert ohnehin.
\phone[mobile]{[+49 151 0000000]}
\email{[vorname.nachname@example.com]}
% ATS: Profil-URLs sichtbar als Text ausgeben. \href{...}{LinkedIn} speichert
% die Adresse in einer PDF-Annotation, die die meisten Parser nie lesen - sie
% sehen nur das Wort "LinkedIn" und extrahieren gar keine Adresse.
\extrainfo{\href{[https://linkedin.com/in/dein-profil]}{linkedin.com/in/[dein-profil]} \textbar{}
           \href{[https://github.com/dein-username]}{github.com/[dein-username]}}

\begin{document}

\makecvtitle

% ============================================================
%     KURZPROFIL
% ============================================================

\vspace{6pt}
\small{[Drei Zeilen, zugeschnitten auf die Zielstelle. Was gewinnt der Arbeitgeber? Beispiel: \glqq Informatikstudent im vierten Semester, Schwerpunkt Datenverarbeitung. Praktische Erfahrung in Python und SQL aus [Projekt]. Sucht eine Werkstudentenstelle, in der [konkreter Beitrag].\grqq{}]}

% ============================================================
%     KERNKOMPETENZEN
% ============================================================

\section{\headingCompetencies}
\vspace{1pt}
\begin{itemize}

\item \textbf{[Kompetenzfeld 1]}: [Konkrete Werkzeuge, Sprachen, Frameworks.]

\item \textbf{[Kompetenzfeld 2]}: [Konkrete Werkzeuge, Sprachen, Frameworks.]

\item \textbf{[Kompetenzfeld 3]}: [Methoden, Fachwissen, Vorgehensweisen.]

\item \textbf{[Kompetenzfeld 4]}: [Werkzeuge, Plattformen, Software.]

\end{itemize}

% ============================================================
%     BERUFSERFAHRUNG
% ============================================================
% Datumsformat MM/JJJJ, Bereichstrennung mit einfachem ASCII-Bindestrich.
% Ein doppeltes -- wird zum Halbgeviertstrich (U+2013); viele Parser
% erkennen dann keinen Zeitraum. Siehe 05-cv-templates.md.

\section{\headingExperience}
\vspace{3pt}

\cventry{[MM/JJJJ-heute]}{[Position]}{[Unternehmen]}{[Ort]}{}{
\begin{itemize}
    \item {[Aufgabe oder Erfolg 1 - konkret, moeglichst mit Zahlen.]}
    \item {[Aufgabe oder Erfolg 2]}
    \item {[Aufgabe oder Erfolg 3]}
\end{itemize}}

\cventry{[MM/JJJJ-MM/JJJJ]}{[Position]}{[Unternehmen]}{[Ort]}{}{
\begin{itemize}
    \item {[Aufgabe oder Erfolg 1]}
    \item {[Aufgabe oder Erfolg 2]}
\end{itemize}}

% ============================================================
%     PROJEKTE
% ============================================================
% Feste Eintragsform, damit generierte Lebenslaeufe das Format nicht
% jedes Mal neu erfinden. Tragende Sektion fuer das Informatik-Profil.

\section{\headingProjects}
\vspace{3pt}

\cventry{[MM/JJJJ]}{[Technologien / Rolle]}{[Projektname]}{}{}{
\begin{itemize}
    \item {[Was es tut und was du gebaut hast - eine Zeile.]}
    \item {[Ergebnis, Umfang oder Wirkung. Beleg verlinken: \href{[URL]}{github.com/[user]/[repo]}]}
\end{itemize}}

% ============================================================
%     AUSBILDUNG
% ============================================================
% Argument 5 ist das Notenfeld. Bei laufendem Studium den
% voraussichtlichen Abschluss angeben - ein geschlossener Zeitraum
% liest sich sonst wie ein bereits erworbener Abschluss.

\section{\headingEducation}
\vspace{3pt}

\cventry{[MM/JJJJ-heute]}{[Studiengang]}{[Hochschule]}{[Ort]}{Note: [X,X]}{
[Schwerpunkte: Themen.] Voraussichtlicher Abschluss: [MM/JJJJ].
}

\cventry{[MM/JJJJ-MM/JJJJ]}{[Abschluss]}{[Schule / Hochschule]}{[Ort]}{Note: [X,X]}{
[Relevante Schwerpunkte oder Abschlussarbeit.]
}

% ============================================================
%     SPRACHEN
% ============================================================

\section{\headingLanguages}
\vspace{1pt}
\begin{itemize}
\item {[Sprache 1] (Muttersprache), [Sprache 2] ([Niveau, z.\,B. C1]), [Sprache 3] ([Niveau]).}
\end{itemize}

% ============================================================
%     AUSZEICHNUNGEN (optional)
% ============================================================
% Im Standardumfang deaktiviert, damit das 1-Seiten-Budget aufgeht.
% Einkommentieren und dafuer an anderer Stelle streichen.
% Hinweis: einfacher ASCII-Bindestrich, kein --. Ein Halbgeviertstrich
% (U+2013) landet unter lualatex als ungueltiges Byte im Textlayer.
%
% \section{\headingAwards}
% \vspace{1pt}
% \begin{itemize}
% \item {\textbf{[Auszeichnung]} - [Veranstaltung / Organisation] ([Jahr]).}
% \end{itemize}

% ============================================================
%     PUBLIKATIONEN - nur akademische Laufbahn (PhD-Track)
% ============================================================
% Fuer Werkstudenten- und Industriestellen bewusst deaktiviert: die Sektion
% ist akademischer Ballast und verleitet zum Auffuellen. Nur einkommentieren,
% wenn die Zielstelle wissenschaftliche Publikationen ausdruecklich erwartet.
%
% \newcommand{\headingPublications}{Publikationen}
% \section{\headingPublications}
% \vspace{1pt}
% \begin{itemize}
% \item {[Autor(en)] ([Jahr]). [Titel]. [Journal/Konferenz]. \href{[DOI_URL]}{DOI}}
% \end{itemize}

% ============================================================
%     KEINE REFERENZEN-SEKTION
% ============================================================
% "Referenzen auf Anfrage" ist eine angloamerikanische Konvention ohne
% Funktion in Deutschland - Arbeitszeugnisse liegen als Anlagen bei.
% Die Zeile kostet eine Ueberschrift und transportiert null Information.

\end{document}
